% !TEX root = Chicken Fajita Sheet Pan Meal Prep.tex
\documentclass[14pt]{extarticle}
\usepackage[utf8]{inputenc}
% --- REMOVE PAGE NUMBERING ---
\pagestyle{empty}
% --- MARGIN SETUP ---
\usepackage[left=0.75in, right=0.75in, top=0.5in, bottom=0.5in]{geometry}
% --- LIST SETTINGS ---
\usepackage{enumitem}
\setlist[itemize]{leftmargin=1em}
% Labels are centered in a fixed-width box so plain bullets line up with the
% wider prep icons (\veg, \measure) below.
\SetLabelAlign{center}{\makebox[\labelwidth][c]{#1}}
\setlist[itemize]{leftmargin=1em, itemsep=0.2em, label={\raisebox{-0.1ex}{\large\textbullet}},
    align=center, labelwidth=0.9em, labelsep=0.4em}
\setlist[enumerate]{leftmargin=1.5em, label=\textbf{\arabic*.}, itemsep=0.25em}
\usepackage{parskip}
% --- CONTROL SECTION FONT SIZE ---
\usepackage{titlesec}
\titleformat{\section}{\large\bfseries}{}{0em}{}
\titleformat{\subsection}{\normalsize\bfseries}{}{0em}{}
\titlespacing*{\subsection}{0pt}{0.75em}{0.25em}
\titlespacing*{\section}{0pt}{1em}{0.5em}
% --- FONT SETUP ---
\usepackage{helvet}
\renewcommand{\familydefault}{\sfdefault}
% --- CUSTOM COMMAND FOR TITLE ---
\newcommand{\makeRecipeTitle}[1]{
    {\noindent \Large \textbf{#1}}
    \par \vspace{-0.75em}
}
% --- CUSTOM COMMAND FOR SERVINGS ---
% Usage: \recipeServings{4--5} or \recipeServings{6 (one burrito each)}
% Place directly after the first \makeRecipeTitle on page 1.
\newcommand{\recipeServings}[1]{
    {\noindent \itshape Serves #1 \par}
    \vspace{-0.75em}
}
% --- PREP ICONS ---
% Usage: \item[\veg] Red Bell Pepper – 1, sliced
%        \item[\measure] Fish Sauce – 1 1/2 tbsp (22ml)
% Replace the bullet on lines to handle before cooking starts: produce to cut
% (carrot), and spices, seasonings, sauces, and pastes with a measured amount
% (measuring spoon, drawn with TikZ). "To taste" lines get neither.
\usepackage{fontawesome5}
\usepackage{tikz}
\newcommand{\veg}{\raisebox{0.25ex}{\small\faCarrot}}
\newcommand{\measure}{\tikz[x=0.9em, y=0.9em,
        baseline={([yshift=-0.36em]current bounding box.center)}]{
    \begin{scope}[rotate=45]
        % Bowl with a rim highlight, and a flat handle with a hanging hole.
        \fill (0,0) ellipse [x radius=0.33, y radius=0.29];
        \draw[white, line width=0.035em, line cap=round]
            (-0.1,0.173) arc[start angle=120, end angle=210, radius=0.2];
        \fill (0.2,-0.07) -- (0.86,-0.105) arc[start angle=-90, end angle=90, radius=0.105]
            -- (0.2,0.07) -- cycle;
        \fill[white] (0.86,0) circle (0.04);
    \end{scope}
}}
% --- CUSTOM COMMAND FOR SOURCE CREDIT ---
% Usage: \recipeSource{Creator -- Video Title}{https://...}
% Place just before \end{document}; pins the credit to the bottom of the last page.
\usepackage[hidelinks]{hyperref}
\urlstyle{same}
\newcommand{\recipeSource}[2]{
    \vfill
    {\noindent \footnotesize \textbf{Source:} #1 \\ \url{#2}}
}


\begin{document}

% ================= PAGE 1 =================
\makeRecipeTitle{Chicken Fajita Sheet Pan Meal Prep}
\recipeServings{5 (meal-prep containers)}

\section*{Ingredients}

\textbf{Fajita Spice Mix}
\begin{itemize}[itemsep=0.1em, label={\raisebox{-0.1ex}{\large\textbullet}}]
      \item[\measure] Chili Powder – 4 tsp (11g)
      \item[\measure] Ground Cumin – 2½ tsp (6g)
      \item[\measure] Smoked Paprika – 2 tsp (5g)
      \item[\measure] Onion Powder – 1 tsp (3g)
      \item[\measure] Dried Oregano – 1 tsp (1g) (Mexican oregano if you have it)
      \item[\measure] Brown Sugar – 1 tsp (4g)
      \item[\measure] Kosher Salt – 2 tsp
      \item[\measure] Black Pepper – ½ tsp
\end{itemize}

\textbf{The Chicken}
\begin{itemize}[itemsep=0.1em, label={\raisebox{-0.1ex}{\large\textbullet}}]
      \item Chicken Thighs – 2.2 lbs (1kg), boneless and skinless
      \item[\cut] Limes – 2, zested and juiced
      \item[\cut] Garlic Cloves – 3, grated
      \item[\measure] Olive Oil – 2 tbsp (30ml)
      \item[\measure] Chipotle Sauce – 1 tbsp (e.g., La Costeña) (optional)
\end{itemize}

\textbf{The Peppers \& Onions}
\begin{itemize}[itemsep=0.1em, label={\raisebox{-0.1ex}{\large\textbullet}}]
      \item[\cut] Bell Peppers – 4 (mixed colors), sliced into ½-inch strips
      \item[\cut] Yellow Onions – 2 medium, sliced into ½-inch strips
      \item[\measure] Olive Oil – 1 tbsp (15ml)
\end{itemize}

\textbf{To Serve}
\begin{itemize}[itemsep=0.1em, label={\raisebox{-0.1ex}{\large\textbullet}}]
      \item \recipeLink{Frozen Cilantro Rice}{Frozen Cilantro Rice} and \recipeLink{Frozen Black Beans}{Frozen Black Beans} (or plain rice)
      \item[\cut] Cilantro – a small handful, chopped
      \item[\cut] Lime – 1, cut into wedges
      \item Salsa, Sour Cream, or Hot Sauce – to taste
\end{itemize}

\newpage
\makeRecipeTitle{Chicken Fajita Sheet Pan Meal Prep}
\section*{Instructions}

\begin{enumerate}

      \item \textbf{Prep Everything} \\
            Slice the bell peppers and onions into strips, grate the garlic, and zest and juice the 2 limes. Chop the cilantro, cut the serving lime into wedges, and measure out everything marked with a spoon.

      \item \textbf{Marinate the Chicken} \\
            Stir the spice mix together and set aside 1½ tablespoons for the vegetables. In a large bowl, whisk the rest with the lime zest and juice, garlic, olive oil, and chipotle sauce. Add the chicken and toss until evenly coated. Marinate for at least 30 minutes, or overnight for deeper flavor, but no longer than 24 hours, or the lime juice makes the chicken mushy.

      \item \textbf{Preheat} \\
            Set one oven rack in the upper third and one in the lower third, and preheat the oven to 425$^\circ$F while the chicken marinates.

      \item \textbf{Roast the Vegetables} \\
            Toss the peppers and onions with the olive oil and the reserved spice mix, and spread them onto a large lined sheet pan (18 $\times$ 13 in) in a single layer. Roast on the upper rack on their own for 10 minutes. Don't stir them yet.

      \item \textbf{Add the Chicken} \\
            Lay the chicken flat on a second lined tray, slide it onto the lower rack, and set a new 10-minute timer. When it goes off, toss the vegetables (their only toss) and return them to the upper rack. Keep cooking until the vegetables are soft and charred at the edges and the chicken reaches 165$^\circ$F (a further 10--15 minutes). \\
            \textbf{Chef's Tip:} For more char, move the chicken to the upper rack and broil for the last 2--3 minutes. Watch it closely.

      \item \textbf{Rest and Slice} \\
            Let the chicken rest for 5 minutes, then slice it into thin strips. Toss the strips with the vegetables and any tray juices.

      \item \textbf{Portion and Store} \\
            Divide the rice and beans between 5 meal-prep containers, top with the fajita chicken and vegetables, and garnish with cilantro. Add a lime wedge to each, and keep the salsa or sour cream on the side until serving. Keeps for 4 days in the fridge; microwave until hot (1--2 minutes), or reheat in a skillet for tacos.

\end{enumerate}

\end{document}
